\documentclass[../../main.tex]{subfiles}

\begin{document}

\chapter[Implementasi Lexical Analyzer]{Implementasi Lexical Analyzer}
\label{chap:implementasi-lexer}

\begin{outcomebox}
\textbf{Sub-CPMK:} \subcpmk{2.4}\\[3pt]
Setelah menyelesaikan bab ini, mahasiswa mampu:
\begin{enumerate}[label=\arabic*.,leftmargin=*]
  \item menjelaskan arsitektur scanner dengan buffer dan lookahead;
  \item mengimplementasikan lexer subset bahasa C versi
        \textit{hand-written} dalam C;
  \item mengimplementasikan lexer versi Flex;
  \item menguji keluaran kedua lexer pada berkas masukan uji;
  \item menghasilkan diagnostik kesalahan leksikal dan
        mengintegrasikan lexer ke proyek Subset-C (M1).
\end{enumerate}
\end{outcomebox}

\section*{Materi Pokok}
\begin{itemize}[leftmargin=*]
  \item Buffer masukan, indeks \code{start} dan \code{current}, serta
        lookahead.
  \item Struktur dan implementasi lexer tangan dalam C.
  \item Lexer generator-based dengan Flex.
  \item Pengujian kedua lexer pada berkas masukan uji.
  \item Penanganan kesalahan, diagnostik, dan integrasi proyek.
\end{itemize}

\subfile{section-01}
\subfile{section-02}
\subfile{section-03}
\subfile{section-04}
\subfile{section-05}

\section*{Contoh}
\addcontentsline{toc}{section}{Contoh}

\begin{contoh}[Menelusuri lexer]
Untuk masukan \code{a <= b;}, lexer memanggil \code{skip\_space},
mencatat posisi 1:1, lalu memindai identifier \code{a}. Token
pertama adalah \code{IDENT} dengan lexeme \code{a}. Setelah spasi
dilewati, lexer membaca \code{<} dan melihat \code{=} melalui
\code{peek}. Karena \code{<=} cocok, lexer memajukan posisi sekali
lagi dan menghasilkan \code{LE}. Berikutnya \code{b} menghasilkan
\code{IDENT}, lalu \code{;} menghasilkan \code{SEMI}. Akhirnya
\code{EOF} dikembalikan pada baris berikutnya.
\end{contoh}

\begin{contoh}[Pemulihan kesalahan]
Untuk masukan \code{int a = 3.;}, lexer menghasilkan token
\code{INT}, \code{IDENT(a)}, \code{ASSIGN}, dan \code{INT\_LIT(3)}.
Ketika melihat titik tanpa digit berikutnya, lexer tidak
memasukkannya ke bilangan. Pada pemanggilan berikutnya, titik
menghasilkan \code{TOK\_ERROR}. Lexer mencatat
\code{baris 1, kolom 10: karakter tidak dikenal}, lalu melanjutkan
ke \code{;}. Parser menerima token kesalahan dan dapat melaporkan
masalah tanpa kehilangan sisa program.
\end{contoh}

\begin{contoh}[Prioritas aturan Flex]
Pada masukan \code{a <= b}, aturan \code{"<="} dan aturan \code{"<"}
sama-sama cocok pada posisi \code{<}. Flex memilih pattern
terpanjang, yaitu \code{"<="}. Jika kedua pattern memiliki panjang
sama, aturan yang lebih awal dalam berkas yang menang. Karena itu,
operator dua karakter harus ditulis sebelum operator satu karakter.
\end{contoh}

\begin{contoh}[Perbandingan ukuran kode]
Lexer tangan pada proyek ini sekitar 350 baris C, sedangkan lexer
Flex sekitar 85 baris. Lexer tangan memberikan kontrol penuh atas
lokasi, pesan kesalahan, dan penanganan kasus khusus. Lexer Flex
lebih ringkas dan mudah diubah, tetapi memerlukan generator dan
sedikit lebih sulit memberikan diagnostik yang sangat spesifik.
\end{contoh}

\begin{contoh}[Membaca keluaran tes]
Berkas \file{basic.expected} memuat baris
\code{1:5 IDENT main}. Artinya, pada baris 1 kolom 5, lexer
menghasilkan token \code{IDENT} dengan lexeme \code{main}. Jika
keluaran aktual berbeda, periksa apakah lexer menghitung kolom
mulai dari 1 dan apakah baris baru menaikkan nomor baris.
\end{contoh}

\section*{Aktivitas Pembelajaran dan Praktikum}
\addcontentsline{toc}{section}{Aktivitas Pembelajaran dan Praktikum}

\begin{table}[H]
\centering
\caption{Rencana aktivitas Bab VII}
\label{tab:ch07-aktivitas}
\begin{tabularx}{\textwidth}{@{}llX@{}}
\toprule
\textbf{Sesi} & \textbf{Bentuk} & \textbf{Kegiatan} \\
\midrule
1 & Ceramah & Pola scanner: \code{start}, \code{current}, dan lookahead \\
2 & Live coding & Lexer tangan: identifier, angka, operator \\
3 & Live coding & Lexer Flex dan perbandingan aturan \\
4 & Praktikum & M1: implementasi, pengujian, dan laporan \\
\bottomrule
\end{tabularx}
\end{table}

\subsection*{M1 / Praktikum 1: Lexer}
\textbf{Tujuan.} Menghasilkan lexer Subset-C yang bekerja, diuji, dan
terdokumentasi.

\textbf{Keluaran.}
\begin{enumerate}[label=\arabic*.,leftmargin=*]
  \item Lexer tangan dalam C (\file{lexer.c}) yang lulus tes.
  \item Lexer Flex (\file{lexer\_flex.l}) yang menghasilkan token
        setara.
  \item Minimal 30 kasus uji dengan berkas harapan.
  \item Keluaran program: aliran token kedua lexer pada berkas
        masukan uji.
  \item Laporan singkat: desain, hasil tes, kendala, dan refleksi.
\end{enumerate}

\textbf{Langkah kerja.}
\begin{enumerate}[label=\arabic*.,leftmargin=*]
  \item Lengkapi \file{include/token.h} dan \file{include/lexer.h}.
  \item Implementasikan \file{src/lexer.c} dan \file{src/error.c}.
  \item Implementasikan \file{src/main\_lex.c}.
  \item Implementasikan \file{src/lexer\_flex.l}.
  \item Jalankan \file{build.bat} dan \file{run-tests.bat}.
  \item Bandingkan keluaran lexer tangan dan Flex.
  \item Tulis laporan dan serahkan.
\end{enumerate}

\textbf{Rubrik.}
\begin{table}[H]
\centering
\caption{Rubrik Praktikum 1}
\label{tab:ch07-rubrik}
\begin{tabularx}{\textwidth}{@{}lXr@{}}
\toprule
\textbf{Aspek} & \textbf{Indikator} & \textbf{Bobot} \\
\midrule
Kebenaran & Semua token dikenali, tes lulus & 40\% \\
Kualitas kode & Struktur, penamaan, komentar, tanpa peringatan & 30\% \\
Kelengkapan & Tangan dan Flex, 30 tes, laporan & 20\% \\
Pengujian & Kasus normal, batas, dan kesalahan & 10\% \\
\bottomrule
\end{tabularx}
\end{table}

\textbf{Definisi selesai.}
\begin{checklist}
  \item Lexer tangan dan Flex berjalan.
  \item Minimal 30 tes lulus.
  \item Keluaran kedua lexer identik pada berkas uji.
  \item Diagnostik menyebut lokasi dan penyebab.
  \item Laporan diserahkan.
\end{checklist}

\section*{Latihan dan Refleksi}
\addcontentsline{toc}{section}{Latihan dan Refleksi}

\subsection*{Latihan}
\begin{enumerate}[label=\arabic*.,leftmargin=*]
  \item Jelaskan fungsi \code{lexeme\_begin} dan \code{forward}.
  \item Mengapa sentinel diperlukan pada buffer?
  \item Tuliskan alur \code{lexer\_next} untuk identifier.
  \item Bagaimana lexer membedakan integer dan float?
  \item Mengapa operator dua karakter harus diperiksa lebih dahulu?
  \item Rancang pemulihan untuk komentar blok yang tidak ditutup.
  \item Bandingkan lexer tangan dan Flex dari sisi kontrol dan
        produktivitas.
  \item Tambahkan dukungan komentar bersarang pada lexer Flex.
\end{enumerate}

\subsection*{Refleksi}
\begin{reflectionbox}
\begin{enumerate}[label=\arabic*.,leftmargin=*]
  \item Bagian implementasi mana yang paling menantang? Mengapa?
  \item Apakah tes Anda mencakup kasus batas yang cukup?
  \item Jika Anda harus memilih satu pendekatan untuk proyek nyata,
        mana yang Anda pilih dan mengapa?
\end{enumerate}
\end{reflectionbox}

\section*{Asesmen}
\addcontentsline{toc}{section}{Asesmen}

Bab VII memuat Praktikum 1 dengan bobot 12\% sesuai RPS. Asesmen
mencakup kebenaran, kualitas kode, kelengkapan, dan pengujian.

\begin{table}[H]
\centering
\caption{Instrumen asesmen Bab VII}
\label{tab:ch07-asesmen}
\begin{tabularx}{\textwidth}{@{}lXl@{}}
\toprule
\textbf{Instrumen} & \textbf{Indikator} & \textbf{Bobot} \\
\midrule
Praktikum 1 & Lexer tangan dan Flex, 30 tes, keluaran uji, laporan &
12\% \\
Kuis & Buffer, maximal munch, pemulihan & formatif \\
Review kode & Keterbacaan, struktur, penanganan kesalahan & bagian P1 \\
\bottomrule
\end{tabularx}
\end{table}

Kriteria kelulusan: lexer mengenali seluruh token Subset-C, lulus
minimal 30 tes, menghasilkan diagnostik yang informatif, dan
terdokumentasi dengan baik.

\section*{Checklist Pencapaian Kompetensi}
\addcontentsline{toc}{section}{Checklist Pencapaian Kompetensi}

\begin{checklist}
  \item Saya dapat menjelaskan skema buffer dan sentinel.
  \item Saya dapat menjelaskan peran lexeme begin dan forward.
  \item Saya dapat mengimplementasikan lexer tangan dalam C.
  \item Saya dapat mengimplementasikan lexer dengan Flex.
  \item Saya dapat menangani kesalahan leksikal tanpa berhenti.
  \item Saya dapat menghasilkan diagnostik dengan lokasi.
  \item Saya dapat menulis tes dan berkas harapan.
  \item Saya dapat menguji kedua lexer pada berkas masukan uji dan
        membandingkan keluarannya.
  \item Saya sudah menyelesaikan M1/Praktikum 1.
\end{checklist}

\section*{Rangkuman}
\addcontentsline{toc}{section}{Rangkuman}

Pada bab ini telah dibahas:
\begin{itemize}[leftmargin=*]
  \item buffer masukan, sentinel, lexeme begin, dan forward pointer;
  \item implementasi lexer tangan dalam C;
  \item implementasi lexer dengan Flex dan perbandingannya;
  \item strategi penanganan dan pelaporan kesalahan leksikal;
  \item driver, pengujian, dan integrasi proyek M1.
\end{itemize}

Bab VIII memulai analisis sintaksis. Kita akan beralih dari token ke
tata bahasa bebas konteks, mempelajari derivasi, pohon urai, dan cara
merancang tata bahasa untuk Subset-C.

\section*{Evaluasi}
\addcontentsline{toc}{section}{Evaluasi}

\begin{enumerate}[label=\arabic*.,leftmargin=*]
  \item Jelaskan pola \code{start} dan \code{current} pada scanner, serta
        mengapa scanner tetap berjalan setelah menemukan kesalahan.
  \item Mengapa lexer memerlukan lookahead? Berikan dua contoh.
  \item Tuliskan pseudokode \code{lexer\_next} untuk identifier.
  \item Jelaskan perbedaan penanganan integer dan float.
  \item Rancang aturan Flex untuk operator dua karakter.
  \item Bagaimana lexer memulihkan diri dari karakter ilegal?
  \item Mengapa format diagnostik \code{berkas:baris:kolom} penting?
  \item Bandingkan lexer tangan dan Flex untuk proyek berukuran
        besar.
  \item Jalankan kedua lexer pada tiga berkas masukan uji dan
        bandingkan aliran token yang dihasilkan.
\end{enumerate}

\subsection*{Kunci dan Rubrik Singkat}
\begin{itemize}[leftmargin=*]
  \item Soal 1--4: ketepatan konsep (70\%), contoh (30\%).
  \item Soal 5--6: kebenaran aturan dan pemulihan (70\%),
        penjelasan (30\%).
  \item Soal 7--9: analisis (60\%), argumentasi (40\%).
\end{itemize}


\ifSubfilesClassLoaded{%
  \bibliographystyle{apalike}%
  \bibliography{../../references}%
}{}

\end{document}
